\section{Actual nonlinear smooth-centre normal layers and canonical scalars}

The original split normal layers already have the canonical quotient-ring
action obtained from ambient multiplication and internally proved ideal
annihilation. We now compare them with the literal power-image layers of
an actual original nonlinear smooth-centre stalk. This is an explicit
local-algebra and chart ingredient for the Bierstone--Milman
Hilbert--Samuel equisingularity route, not a sheaf normal-flatness theorem
assumed from that source. A genuine regular-centre flattening supplies
the coordinate equivalence and the WHOLE original vanishing ideal.

\begin{lemma}[Whole original power quotients under a genuine ring equivalence]
\label{hir:nonlinear-centre-actual-power-quotient}
\source{bierstone-milman08-functoriality}
\uses{hir:relative-normal-actual-layer-centre-scalars}
\lean{Hironaka.HironakaSmoothCentreCanonicalLayers.centreQuotientEquiv,Hironaka.HironakaSmoothCentreCanonicalLayers.centreQuotientEquiv_mk,Hironaka.HironakaSmoothCentreCanonicalLayers.centreQuotientEquiv_symm_mk,Hironaka.HironakaSmoothCentreCanonicalLayers.powerLayerQuotientEquiv,Hironaka.HironakaSmoothCentreCanonicalLayers.powerLayerQuotientEquiv_mk,Hironaka.HironakaSmoothCentreCanonicalLayers.powerLayerQuotientEquiv_symm_mk}
\leanok
Let $E:R\simeq S$ be a genuine commutative ring equivalence and let
$I,J\subset R$. For every $q$, the original quotient map induces a
genuine equivalence
\[
 R/(I+J^{q+1})\simeq
 S/(E(I)+E(J)^{q+1}),\qquad [g]\longmapsto[E(g)].
\]
It also gives the centre-ring equivalence $R/J\simeq S/E(J)$ with the
same actual representative formula. Both inverse maps send $[h]$ to
$[E^{-1}(h)]$ and retain entire original ideals, not a finite selection
of relations.
\end{lemma}
\begin{proof}
The actual image map preserves ideal sums and products, so it sends
$I+J^{q+1}$ to $E(I)+E(J)^{q+1}$. Its inverse has the inverse ideal
dictionary because $E$ is bijective. Thus both actual ring maps descend
to the full quotient ideals. Their representative compositions are the
identity, so quotient surjectivity proves the full inverse identities.
The same construction with $J$ proves the centre-ring statement.
\end{proof}

\begin{lemma}[Every original power-image layer has its actual nonlinear comparison]
\label{hir:nonlinear-centre-actual-layer-comparison}
\source{bierstone-milman08-functoriality}
\uses{hir:nonlinear-centre-actual-power-quotient,hir:relative-normal-graded-literal-layers}
\lean{Hironaka.HironakaSmoothCentreCanonicalLayers.actualPowerLayerAddEquiv,Hironaka.HironakaSmoothCentreCanonicalLayers.actualPowerLayerAddEquiv_val,Hironaka.HironakaSmoothCentreCanonicalLayers.actualPowerLayerAddEquiv_symm_val,Hironaka.HironakaSmoothCentreCanonicalLayers.actualPowerLayerAddEquiv_class}
\leanok
For every $q$, the preceding whole quotient equivalence restricts to an
additive equivalence between the literal images
\[
 L_q=\operatorname{image}(J^q\to R/(I+J^{q+1})),\quad
 L'_q=\operatorname{image}(E(J)^q\to S/(E(I)+E(J)^{q+1})).
\]
EVERY actual represented class $[g]_q$, $g\in J^q$, is sent to
$[E(g)]_q$; the inverse has the corresponding $E^{-1}$ formula. No
desired layer map, surjectivity or comparison equality is an input.
\end{lemma}
\begin{proof}
Take the actual quotient equivalence on each original layer element.
Its genuine power representative $g$ belongs to $J^q$, so $E(g)$ belongs
to $E(J)^q$ and its class lies in the target image. Conversely every
target power element has an actual preimage in $J^q$ because ideal
images under the genuine bijection are precisely set images, and images
commute with powers. This proves closure for the true inverse quotient
map. Restricting the whole quotient inverse identities proves the
restricted maps inverse; addition follows from their original ring maps.
All class formulas hold on every representative because they are the
actual quotient representative formulas, independently of choices.
\end{proof}

\begin{lemma}[Naturality of the intrinsically descended centre actions]
\label{hir:nonlinear-centre-canonical-scalars-and-whole-sum}
\source{bierstone-milman08-functoriality}
\uses{hir:nonlinear-centre-actual-layer-comparison,hir:relative-normal-actual-associated-centre-scalars}
\lean{Hironaka.HironakaSmoothCentreCanonicalLayers.actualPowerLayerAddEquiv_centre_smul,Hironaka.HironakaSmoothCentreCanonicalLayers.actualPowerLayerCentreEquiv,Hironaka.HironakaSmoothCentreCanonicalLayers.actualPowerAssociatedAddEquiv,Hironaka.HironakaSmoothCentreCanonicalLayers.actualPowerAssociatedAddEquiv_component,Hironaka.HironakaSmoothCentreCanonicalLayers.actualPowerAssociatedAddEquiv_of,Hironaka.HironakaSmoothCentreCanonicalLayers.actualPowerAssociatedAddEquiv_centre_smul,Hironaka.HironakaSmoothCentreCanonicalLayers.actualPowerAssociatedCentreEquiv}
\leanok
The actual layer equivalence is semilinear for the actual centre-ring
equivalence. For EVERY centre scalar $c$ and actual layer element $x$,
\[
 e_q(cx)=\overline E(c)e_q(x).
\]
These true degree maps also induce the whole direct-sum semilinear
equivalence, preserving every actual degree inclusion and every actual
power class. All actions on both sides are their independently defined
ambient-product quotient descents.
\end{lemma}
\begin{proof}
Represent $c$ by $r\in R$. The canonical quotient scalar on either side
acts by the genuine ambient product. The underlying whole quotient ring
map sends $rx$ to $E(r)e_q(x)$ because it preserves multiplication.
The centre quotient sends the scalar representative $r$ to $E(r)$,
proving the formula. Quotient scalar surjectivity gives EVERY $c$; the
whole layer carrier, not selected generators, is used. Apply these actual
additive equivalences componentwise to finite-support elements. Their
inverse maps preserve zero, hence finite support, and the inverse laws
hold in every component. The scalar formula holds componentwise, giving
whole semilinearity and the stated degree inclusion identities.
\end{proof}

\begin{theorem}[Original regular-centre data supply the whole nonlinear stalk bridge]
\label{hir:nonlinear-centre-original-regular-stalk-bridge}
\source{bierstone-milman08-functoriality}
\uses{hir:nonlinear-centre-canonical-scalars-and-whole-sum,hir:smooth-centre-entire-vanishing-stalk-normal-ideal}
\lean{Hironaka.HironakaSmoothCentreCanonicalLayers.wholeChartStalkIdeal_le_of_original_le,Hironaka.HironakaSmoothCentreCanonicalLayers.exists_original_regular_canonical_layer_bridge}
\leanok
Let $C$ be an actual original centre regular of dimension $d$ at $x$ in
an original complex manifold of dimension $d+t$. Let $A$ be an original
ideal and assume its actual germ-ideal inclusion in the vanishing ideal
of $C$. Put $I_x$ equal to its WHOLE original chart stalk and $K_x$ equal
to the WHOLE original vanishing chart stalk of $C$. Original regularity
internally supplies an actual nonlinear flattening with genuine chart
domains, a germ equivalence $E$, and $K_x=E(J)$. It supplies
$I_x\subset K_x$, the actual centre-ring comparison and all original
split layers and their whole direct sum for $E^{-1}(I_x)$, compared
semilinearly to the literal original $K_x$-power image layers of $I_x$.
Every represented germ and every degree inclusion is retained by the
true $E$ formula. No coordinate, ideal, power, layer, scalar or flatness
comparison certificate is supplied by the caller.
\end{theorem}
\begin{proof}
Apply the genuine original regular-centre producer. It constructs an
actual holomorphic local equivalence, containing the actual chart point
in its source, mapping the original centre exactly onto the coordinate
plane on its genuine domain, and fixing zero after recentering. Its
entire induced germ equivalence identifies the WHOLE vanishing stalk
with the image of the true coordinate normal ideal. The actual original
germ-ideal inclusion gives inclusion of every original membership germ;
spanning all such germs proves $I_x\subset K_x$. Bijectivity gives
$E(E^{-1}(I_x))=I_x$ internally. The preceding whole quotient, layer and
canonical scalar constructions applied to this actual $E$ and these
actual whole ideals give every required comparison. The coordinate
source/target and centre formula are retained from the producer, rather
than replaced by an abstract desired chart dictionary.
\end{proof}

This bridge does not construct the nonlinear coherent family on an
entire centre neighbourhood, the intrinsic sheaf module or its stalk
identification. In particular it does not supply an unproved flatness
assumption at the original centre point. Original-input freeness on a
genuine nearby geometric centre domain, chart independence, ambient
mixed-degree Hilbert--Samuel comparison, blowing-up descent and the
singular-centre driver remain further producers. Higher whole truncations
are not assigned a centre quotient-ring action.
